\SourcePage{565}{568}
\section{Spin in a time-dependent magnetic field}

Consider an electrically neutral particle with a magnetic moment, placed in a
uniform magnetic field which varies with time. The particle may be elementary
(a neutron, for example) or composite (an atom). We suppose the field to be so
weak that the particle's magnetic energy in it is small compared with the
separations between its energy levels. The motion of the particle as a whole
may then be considered for a specified internal state.

Let $\hat{\mathbf s}$ denote the operator of the particle's ``intrinsic''
angular momentum: its spin for an elementary particle, or the total angular
momentum $\mathbf J$ for an atom. Write the magnetic-moment operator in the
form (111.1). The Hamiltonian governing the motion of the neutral particle as
a whole is
\begin{equation}
 \widehat H=-\frac{\mu}{s}\hat{\mathbf s}\mathbf H
 \tag{114.1}
\end{equation}
(only the spin-dependent part of the Hamiltonian has been displayed).

In a uniform field this operator contains no explicit coordinates\footnote{The
same reasoning also applies when a particle (including a charged one) moves
through a nonuniform magnetic field and its motion can be treated
semiclassically. The magnetic field encountered as the particle advances
along its trajectory may then simply be regarded as a function of time, and
the same equations may be used for the change of the spin wave function.}.
The particle's wave function therefore separates into coordinate and spin
factors. The former is merely the wave function for free motion; below we are
concerned only with the spin factor. We shall show that the problem for a
particle of arbitrary angular momentum $s$ reduces to the simpler problem of
the motion of a spin-$1/2$ particle (E.~Majorana). For this purpose we use the
device already employed in §~57. In place of one particle with spin $s$, we
formally introduce a system of $2s$ spin-$1/2$ ``particles.'' The operator
$\hat{\mathbf s}$ is then represented by the sum
$\sum_a\hat{\mathbf s}_a$ of their spin operators, while the wave function is
represented by a product of $2s$ spinors of rank one. The Hamiltonian (114.1)
accordingly separates into a sum of $2s$ independent Hamiltonians:
\begin{equation}
 \widehat H=\sum_a\widehat H_a,\qquad
 \widehat H_a=-\frac{\mu}{s}\mathbf H\hat{\mathbf s}_a,
 \tag{114.2}
\end{equation}

\SourcePage{566}{569}
so that each of the $2s$ ``particles'' moves independently of all the others.
Once this has been done, the products of components of the $2s$ rank-one
spinors need only be replaced again by the components of an arbitrary
symmetric spinor of rank $2s$.

\begin{center}\textbf{Problems}\end{center}

\ProblemHeading{1.} Find the spin wave function of a neutral spin-$1/2$
particle in a uniform magnetic field whose direction is fixed but whose
magnitude follows an arbitrary function $H(t)$.

\SolutionHeading The wave function is a spinor $\psi^\nu$ satisfying the wave
equation
\begin{equation}
 i\hbar\dot\psi^\nu=-2\mu\mathbf H\hat{\mathbf s}\psi^\nu.
 \tag{1}
\end{equation}
Choose the field direction as the $z$ axis. In spinor components, this equation
then becomes
\[
 i\hbar\dot\psi^1=-\mu H\psi^1,\qquad
 i\hbar\dot\psi^2=\mu H\psi^2.
\]
It follows that
\[
 \psi^1=c_1\exp\left(\frac{i\mu}{\hbar}\int H\,dt\right),
 \qquad
 \psi^2=c_2\exp\left(-\frac{i\mu}{\hbar}\int H\,dt\right).
\]
The constants $c_1$, $c_2$ are fixed by the initial conditions and the
normalization condition $|\psi^1|^2+|\psi^2|^2=1$.

\ProblemHeading{2.} Solve the same problem for a uniform magnetic field of
constant magnitude whose direction rotates uniformly (with angular velocity
$\omega$) about the $z$ axis, maintaining an angle $\theta$ with it.

\SolutionHeading The components of the magnetic field are
\[
 H_x=H\sin\theta\cos\omega t,\qquad
 H_y=H\sin\theta\sin\omega t,\qquad
 H_z=H\cos\theta,
\]
and (1) gives the system
\[
 \dot\psi^1=i\omega_H(\cos\theta\cdot\psi^1
  +\sin\theta\cdot e^{-i\omega t}\psi^2),
\]
\[
 \dot\psi^2=i\omega_H(\sin\theta\cdot e^{i\omega t}\psi^1
  -\cos\theta\cdot\psi^2),
\]
where $\omega_H=\mu H/\hbar$. The substitution
\[
 \psi^1=e^{-i\omega t/2}\varphi^1,\qquad
 \psi^2=e^{i\omega t/2}\varphi^2
\]
reduces these equations to a system of linear equations with constant
coefficients. Solving it gives
\[
 \psi^1=e^{-i\omega t/2}
   \left(c_1e^{i\Omega t/2}+c_2e^{-i\Omega t/2}\right),
\]
\[
 \psi^2=2\omega_H\sin\theta\,e^{i\omega t/2}
 \left[
  \frac{c_1}{\Omega+\omega+2\omega_H\cos\theta}e^{i\Omega t/2}
  -\frac{c_2}{\Omega-\omega-2\omega_H\cos\theta}e^{-i\Omega t/2}
 \right],
\]
where
\[
 \Omega=\sqrt{(\omega+2\omega_H\cos\theta)^2
              +4\omega_H^2\sin^2\theta}.
\]

\SourcePage{567}{570}
